\documentclass[12pt,a4paper]{article}

% Paquetes para español
\usepackage[utf8]{inputenc}
\usepackage[T1]{fontenc}
% Babel with Spanish support would go here if available

% Paquetes adicionales
\usepackage{geometry}
\usepackage{setspace}
\usepackage{csquotes}
\usepackage{hyperref}

% Configuración de márgenes
\geometry{
    left=2.5cm,
    right=2.5cm,
    top=2.5cm,
    bottom=2.5cm
}

% Espaciado
\setstretch{1.15}

% Información del documento
\title{\textbf{LIBROS HERÉTICOS, PROHIBIDOS Y CONTROVERTIDOS EN LA BIBLIOTECA DE FERNÁN DE VELASCO} \\ 
\large Análisis de una colección bibliográfica en la frontera de la ortodoxia (1791)}
\author{}
\date{}

\begin{document}

\maketitle

\section*{INTRODUCCIÓN}

El catálogo de la biblioteca de D. Fernán de Velasco, tasado en 1791, revela un aspecto particularmente intrigante: la presencia de obras incluidas en el Índice de Libros Prohibidos de la Inquisición española, así como textos de autores protestantes, herejes declarados y filósofos ilustrados considerados peligrosos para la fe católica. Esta contradicción aparente ---un miembro del Consejo y Cámara de Su Majestad poseyendo literatura heterodoxa--- requiere un análisis cuidadoso del contexto intelectual, político y religioso de la España del siglo XVIII.

\section{EL ÍNDICE DE LIBROS PROHIBIDOS EN LA BIBLIOTECA}

El propio catálogo registra dos copias del Índice Inquisitorial, lo que demuestra que Velasco conocía perfectamente qué obras estaban prohibidas:

\begin{quote}
\textbf{``Index Librorum prohibitorum a Benedicto XIV 4º Romae 1758 Nº112 12''}
\end{quote}

\begin{quote}
\textbf{``Index delos libros prohibidos en España hasta el año 1767 fol. 2 tom.''}
\end{quote}

La presencia del Índice no es casual. Los eruditos, funcionarios y bibliófilos de alto rango a menudo poseían estas listas precisamente para poder identificar ---y, paradójicamente, buscar--- las obras censuradas que consideraban valiosas para sus investigaciones. El Índice funcionaba así como un catálogo inverso: una guía de lecturas peligrosas pero intelectualmente estimulantes.

\section{VOLTAIRE: EL FILÓSOFO ILUSTRADO PROHIBIDO}

La presencia de Voltaire en la biblioteca Velasco es particularmente significativa. Voltaire (1694-1778) representaba todo lo que la Iglesia Católica combatía: deísmo, anticlericalismo, crítica mordaz de la superstición y defensa de la tolerancia religiosa. Sus obras estaban absolutamente prohibidas en España.

\subsection*{Obras de Voltaire identificadas:}

\begin{quote}
``Voltaire (M. de) Lettres de quelques juifs Portugais, Allemands, et Polonois: 8º 3 tom. 1781 ala Rustica Paris v. 173 42½''
\end{quote}

Esta obra contiene críticas devastadoras al cristianismo, presentadas en forma epistolar satírica.

\begin{quote}
``Voltaire (M. de) Id. El Oráculo de los nuevos Filósofos traducido al Español p. el P. Fr. Pedro Moro 1769 4.º 2 tom. cuad. n. 79 32½''
\end{quote}

Lo más notable es que esta obra está \textbf{traducida al español por un fraile} (P. Fr. Pedro Moro). Esto sugiere una estrategia de refutación: traducir a Voltaire para poder combatirlo mejor. Sin embargo, el mero hecho de poseer estas obras ---incluso en versión ``refutativa''--- era sospechoso.

\textbf{Precio significativo:} 42½ reales por los tres tomos de las \textit{Lettres} indica que Velasco valoraba esta obra, pese a su carácter prohibido.

\section{OBRAS ANTI-PROTESTANTES Y ANTI-HERÉTICAS: LA PARADOJA DEFENSIVA}

La biblioteca contiene numerosas obras de controversia anti-herética, lo que parece ortodoxo. Sin embargo, estas obras citan extensamente a los herejes que refutan, convirtiéndose así en fuentes indirectas de doctrina heterodoxa:

\begin{quote}
\textbf{``Inquisitionibus Supremae Censura Errorum quibus Hæretici Sacras Scripturas arguere dixunt: 4.º Pintia 1554. N.º 78. 10''}
\end{quote}

Obra que cataloga los errores heréticos sobre las Escrituras. Paradójicamente, para refutar los errores, primero había que exponerlos.

\begin{quote}
\textbf{``Ribera (Alonso) Historia Sacra del S. Sacramento contra los Herejes \& estos tiempos 1626 37 24''}
\end{quote}

Historia apologética contra los protestantes que negaban la transubstanciación.

\begin{quote}
\textbf{``Les Artifices des Heretiques 12.º Part. Paris 1681. N. 251. 8''}
\end{quote}

Texto francés sobre las ``artimañas de los herejes''. La lectura de estrategias heréticas podía resultar peligrosamente persuasiva.

\begin{quote}
\textbf{``Henrici VIII Angliae Regis Assertio Septem Sacramentorum adversus Lutherum 16º Parisii 1562 251 4''}
\end{quote}

Obra del rey Enrique VIII de Inglaterra defendiendo los siete sacramentos \textbf{contra Lutero}. Irónico, dado que Enrique VIII posteriormente rompió con Roma. Esta obra, al citar extensamente a Lutero para refutarlo, transmitía doctrinas luteranas al lector.

\section{OBRAS SOBRE JURISDICCIÓN ECLESIÁSTICA: REGALISMO Y TENSIÓN CON ROMA}

Varios textos en la biblioteca abordan el conflictivo tema de la relación entre poder civil y eclesiástico, una cuestión candente en la España borbónica ilustrada:

\begin{quote}
\textbf{``Bosuet (Jacobo Benigno) De Potestate Ecclesiastica 4º Pasta Luxemburgi 1730 279 45''}
\end{quote}

Jacques-Bénigne Bossuet (1627-1704), obispo francés y defensor del galicanismo (doctrina que limitaba el poder papal en Francia). Aunque católico, sus ideas sobre la autonomía de las iglesias nacionales frente a Roma eran vistas con recelo por la Santa Sede.

\begin{quote}
\textbf{``Bosuet (Jacobo Benigno) Histoire des variations des Églises Protestantes 2 tom. Paris 1734 51 48''}
\end{quote}

Historia de las variaciones doctrinales protestantes, obra apologética pero que exponía detalladamente las doctrinas reformadas.

\subsection*{Otras obras de Bossuet presentes:}
\begin{itemize}
    \item ``Exposición de la Doctrina de la Iglesia Católica sobre los matices de las controversias 8º Paris 1748 318 8''
    \item ``Discurso sobre la Historia Universal traducido al Castellano por Salcedo 12º Part. 2 tom. Valencia 1772 172 20''
    \item ``Política Sagrada 4º 3 tom. Mad. 1778. N 12. 96''
    \item ``Variaciones a las Iglesias Protestantes 4º 8 tom. Mad 1772. Nº12. 55''
\end{itemize}

La abundancia de Bossuet (al menos 8 obras diferentes) sugiere simpatías galicanas o regalistas en Velasco.

\begin{quote}
\textbf{``Respuesta de la Santa Sede al exanifiento de Portugal en Italiano Venetia 1760. num 100 pa. 2.º 6''}
\end{quote}

Documento sobre conflictos jurisdiccionales entre Portugal y Roma.

\section{EL CARDENAL TOURNON Y LOS RITOS CHINOS: CONTROVERSIA MISIONERA}

\begin{quote}
\textbf{``Toximon (el Cardenal) Variæ Apologias suyas sobre los Ritos Chineses. 4.º Turino 1710. Part. 2.º numero 100. 8r''}
\end{quote}

Carlos Tomás Maillard de Tournon (1668-1710), legado papal en China, protagonizó la ``Controversia de los Ritos Chinos'': ¿era lícito que los cristianos chinos participaran en ceremonias confucianas de veneración a los ancestros? Los jesuitas decían que sí (eran ritos civiles); otros consideraban que era idolatría.

La presencia de estas apologías indica el interés de Velasco en debates teológicos complejos y geográficamente lejanos, típico de la mentalidad ilustrada cosmopolita.

\section{TRATADOS DE TOLERANCIA Y LIBERTAD RELIGIOSA: LÍMITES PELIGROSOS}

\begin{quote}
\textbf{``L'Esprit des Nations. 12.º 2 volum. 4. Haye 1753. N. 183. 24''}
\end{quote}

Obra francesa sobre el espíritu de las naciones, probablemente relacionada con las ideas ilustradas sobre diversidad cultural y religiosa.

\begin{quote}
\textbf{``Respuesta teológica acerca del abuso de los Escotados a Santiago 1673. W. 258. 8''}
\end{quote}

Debate teológico sobre prácticas consideradas abusivas, mostrando controversias internas católicas.

\section{AUTORES PROTESTANTES O SOSPECHOSOS DE HEREJÍA}

\subsection*{Escritores protestantes en latín:}

La biblioteca contiene obras de autores del norte de Europa (Holanda, Alemania, Inglaterra) en latín, donde la frontera confesional era menos visible. Muchos de estos eruditos eran protestantes:

\begin{quote}
\textbf{``Boxhornii (Marci Zirœii) Institutiones Politicae. 4.º Regio monte 1678. N. 67 15''}
\end{quote}

Marcus Zuerius Boxhorn (1612-1653), filólogo e historiador holandés calvinista. Sus obras sobre política e historia antigua eran valiosas pero provenían de autor protestante.

\begin{quote}
\textbf{``Vossii - (Isaaci) De vera etate Christi: 4.° Haga Comitis 1659 Num. 280 6½''}
\end{quote}

Isaac Vossius (1618-1689), erudito holandés de familia protestante, escribió sobre cronología bíblica. La Haya era centro del protestantismo holandés.

\subsection*{Múltiples obras de Gerardus Joannes Vossius (padre de Isaac):}
\begin{itemize}
    \item Etymologicum Linguae Latinae
    \item De scientiis mathematicis
    \item Ars Historica
    \item Historiae de controversiis Pelagianis
    \item Rhetorices contractae
\end{itemize}

Vossius padre era un erudito protestante liberal cuyos trabajos filológicos e históricos eran admirados internacionalmente, pero su protestantismo era evidente.

\section{OBRAS SOBRE LA REFORMA Y CISMAS RELIGIOSOS}

\begin{quote}
\textbf{``Bosio (Jacomo) Historia della Sacra Religione di Malta fol. 3 tom. 2 volum. Rom. 1594 119 120''}
\end{quote}

Historia de la Orden de Malta, incluye referencias a conflictos con protestantes.

\begin{quote}
\textbf{``Ribera (Francisco) Vida de N.ª Theresia de Jesus 1 1590 172 15''}
\end{quote}

Santa Teresa de Ávila, figura de la Contrarreforma católica, pero su vida estaba llena de experiencias místicas que la Inquisición había investigado con sospecha.

\section{FILOSOFÍA HETERODOXA}

\begin{quote}
\textbf{``Hill (Nicolai) Philosophia Epicurea Democritiana et theophrastica 12.º Parisiis 1604. Atad. 17. 4''}
\end{quote}

Nicholas Hill (c.1570-c.1610), filósofo inglés, defendió una filosofía atomista epicúrea, rechazada por la Iglesia por su materialismo implícito.

\begin{quote}
\textbf{``Sánchez (Francisco) Tractatus Philosophi quod nihil scitur 12.ª Rotterdam 1692 239 4r''}
\end{quote}

Francisco Sánchez (1551-1623), filósofo escéptico portugués-español. Su obra principal \textit{Quod nihil scitur} (Que nada se sabe) cuestionaba radicalmente la posibilidad del conocimiento cierto, incluyendo el teológico. Una obra profundamente subversiva.

\section{EL CONTEXTO: ¿CÓMO POSEÍA VELASCO ESTOS LIBROS?}

Varias explicaciones plausibles:

\subsection*{a) Licencias especiales (\textit{Licentia legendi})}

Los funcionarios de alto rango, teólogos y eruditos podían obtener de la Inquisición permisos especiales para leer libros prohibidos con fines de estudio o refutación. Velasco, como miembro del Consejo Real, probablemente tenía esta licencia.

\subsection*{b) Compra en el extranjero}

Durante viajes diplomáticos o mediante intermediarios, era posible adquirir libros prohibidos en Francia, Italia o los Países Bajos. Las obras de Voltaire compradas en París, las de Bossuet en Luxemburgo, etc., sugieren esta vía.

\subsection*{c) Comercio clandestino}

Existía un mercado negro del libro prohibido. Los libreros discretos podían conseguir obras censuradas para clientes solventes y poderosos.

\subsection*{d) Herencia y donaciones}

Muchos libros prohibidos entraban en bibliotecas por herencia de colecciones anteriores, cuando las obras aún no estaban censuradas.

\subsection*{e) Actitud ilustrada:}

La Ilustración española, incluso en sus versiones ``católicas'', se caracterizaba por una curiosidad intelectual que trascendía las barreras confesionales. Leer a Voltaire no significaba necesariamente compartir sus ideas, sino mantenerse al día en los debates intelectuales europeos.

\section{RIESGOS Y CONSECUENCIAS}

Poseer libros prohibidos no era trivial. Los riesgos incluían:

\begin{itemize}
    \item \textbf{Denuncia ante la Inquisición} (aunque menos frecuente para nobles y funcionarios)
    \item \textbf{Escándalo social} si se hacía público
    \item \textbf{Confiscación de la biblioteca} en casos extremos
    \item \textbf{Daño a la reputación y carrera}
\end{itemize}

Sin embargo, en 1791 la Inquisición española estaba en declive. Carlos III (1759-1788) había limitado severamente sus poderes. Carlos IV (1788-1808) continuó esta política. La Ilustración había penetrado en la administración borbónica.

El hecho de que Antonio Baylo tasara estos libros \textbf{abiertamente}, incluyendo las obras de Voltaire en el catálogo oficial sin censurarlas, indica que en 1791 la posesión de literatura heterodoxa por parte de un alto funcionario ya no era perseguida activamente.

\section{CONCLUSIONES}

La biblioteca de Fernán de Velasco refleja la tensión fundamental de la Ilustración española:

\subsection*{1. Ortodoxia oficial vs. curiosidad intelectual}
\begin{itemize}
    \item Velasco era funcionario real y presumiblemente católico ortodoxo
    \item Pero su biblioteca incluye a Voltaire, Bossuet, filósofos escépticos, eruditos protestantes
\end{itemize}

\subsection*{2. El Índice como catálogo de lectura}
\begin{itemize}
    \item Poseer el Índice de Libros Prohibidos le permitía identificar obras intelectualmente importantes
    \item Los libros prohibidos eran precisamente los más debatidos en Europa
\end{itemize}

\subsection*{3. Regalismo vs. ultramontanismo}
\begin{itemize}
    \item La abundancia de Bossuet sugiere simpatías galicanas/regalistas
    \item Interés en limitar el poder papal a favor de la autoridad real sobre la Iglesia nacional
\end{itemize}

\subsection*{4. Red de conocimiento transnacional}
\begin{itemize}
    \item Compras en París, Lyon, Ámsterdam, La Haya
    \item Velasco participaba en las redes intelectuales europeas ilustradas
\end{itemize}

\subsection*{5. Tolerancia de facto en 1791}
\begin{itemize}
    \item El tasador registró abiertamente obras prohibidas
    \item Indica relajación de la censura en los años finales del Antiguo Régimen
\end{itemize}

\subsection*{6. Funcionalidad práctica de la heterodoxia}
\begin{itemize}
    \item Para refutar ideas peligrosas, había que conocerlas
    \item La literatura anti-herética requería familiaridad con la herejía
    \item Los juristas necesitaban conocer argumentos de todas las partes
\end{itemize}

\subsection*{7. El precio de la transgresión}
\begin{itemize}
    \item Las obras de Voltaire tienen precios medios-altos (32½ y 42½ reales)
    \item No son excesivamente caras, lo que indica que el mercado de libros prohibidos era relativamente accesible para las élites
\end{itemize}

Esta biblioteca documenta así un momento crucial: el ocaso del control inquisitorial absoluto sobre el pensamiento, y el amanecer de una cultura ilustrada española que, sin abandonar formalmente el catolicismo, se abría selectivamente a las corrientes intelectuales europeas, incluyendo aquellas condenadas por la Iglesia.

La paradoja de Velasco ---funcionario real ortodoxo y lector de Voltaire--- es en realidad la paradoja de toda la Ilustración católica: ¿cómo modernizar España intelectualmente sin romper con su identidad religiosa tradicional? Su biblioteca, con sus Voltaires junto a sus tratados anti-heréticos, sus Bossuets junto a sus devocionarios, encarna este dilema histórico.

\section*{ANEXO: LISTA DE OBRAS PROHIBIDAS O CONTROVERTIDAS IDENTIFICADAS}

\subsection*{Autores explícitamente en el Índice:}
\begin{itemize}
    \item Voltaire (\textit{Lettres}, \textit{Oráculo de los nuevos Filósofos})
\end{itemize}

\subsection*{Autores protestantes o heterodoxos:}
\begin{itemize}
    \item Gerardus Joannes Vossius (protestante holandés) - 8 obras
    \item Isaac Vossius (protestante holandés) - 1 obra
    \item Marcus Zuerius Boxhorn (calvinista holandés) - 2 obras
    \item Nicholas Hill (filósofo atomista inglés) - 1 obra
    \item Francisco Sánchez (escéptico radical) - 1 obra
\end{itemize}

\subsection*{Autores galicanos/regalistas:}
\begin{itemize}
    \item Jacques-Bénigne Bossuet - al menos 8 obras
\end{itemize}

\subsection*{Obras sobre herejías (que citan doctrinas heréticas):}
\begin{itemize}
    \item Múltiples tratados anti-luteranos, anti-calvinistas, anti-heréticos
\end{itemize}

\subsection*{Obras sobre controversias jurisdiccionales Iglesia-Estado:}
\begin{itemize}
    \item Tratados sobre poder papal, jurisdicción eclesiástica, conflictos Roma-reinos católicos
\end{itemize}

\subsection*{Obras sobre controversias misioneras:}
\begin{itemize}
    \item Cardenal Tournon sobre Ritos Chinos
\end{itemize}

\textbf{Total estimado de obras potencialmente problemáticas:} Al menos 30-40 obras de las varios miles catalogadas, lo que representa aproximadamente el 1-2\% de la colección. Una proporción pequeña pero significativa, que indica lectura selectiva y consciente de literatura heterodoxa, no mera acumulación casual.

\end{document}
